\subsection{What entire-function growth does and does not specify}
\label{audit:entire-data}

\paragraph{The historical question and its limits.}
After the polynomial formula for \(\chi\), P000568 asks
\ru{Пусть \(R\) и \(P\) --- целые ф-и 1-го пор. роста. Тогда им соотв. потенциалы \(u(x)\). Какие?}
and adds \ru{не обяз. периодических}. The scan confirms the first-order-growth
wording. The strongest supported historical question is to identify the potentials
represented by an entire-function replacement of the finite-gap formula. The
notebook supplies no normalization, admissible zeros, realness assumption, or
global regularity requirement. It therefore does not determine a unique global
inverse problem. This subsection is a verification annotation to that Class F
item, not a declaration that an unformulated open problem has been solved.

\paragraph{Precisely relevant later results.}
\citet[Thm.~5.1, pp.~535--536]{marchenko1975characterization} characterize
periodic and antiperiodic spectra for a real periodic Hill operator in specified
Sobolev classes using a normalized slit-domain conformal map and weighted
conditions on its slit heights. Their paper was received on 3 February 1975.
That theorem has a specified period, reality, normalization and spectral
asymptotics, absent from the seminar question. \citet[Hypotheses~3.1--3.2,
Eq.~(3.14), Lemmas~3.3--3.5]{gesztesy1999infinite}, first circulated on
7 July 1999, develop an infinite-genus construction with convergent spectral
products, an assumed Green-function identity and controlled high-energy
asymptotics. These are concrete realizations of the direction, not an exact
match to arbitrary entire functions of order one. No evidence found in these
sources establishes why Gelfand did not specify or pursue this formulation.

\paragraph{A bounded reconstruction that can be settled.}
Let \(I\subset\mathbb R\) be a nonempty open interval and \(U\subset\mathbb C\)
a connected open set. Suppose \(P(E,x)\ne0\) on \(U\times I\), is twice
continuously differentiable in \(x\), and is holomorphic in \(E\), with its first
two \(x\)-derivatives continuous. Let \(R\) be holomorphic on \(U\), and suppose
there is a holomorphic \(y\) on \(U\) such that \(y^2=R\). For example one can
choose a simply connected region avoiding the zeros of \(R\). All conclusions
below are local in this specified region; entire functions are a special case.

\begin{proposition}[Local compatibility criterion \ours]
\label{prop:entire-compatibility}
Define
\[
 Q(E,x)=E-\frac{R(E)}{P(E,x)^2}
       +\frac{P_{xx}(E,x)}{2P(E,x)}
       -\frac{P_x(E,x)^2}{4P(E,x)^2}.
\]
There is a continuous potential \(u:I\to\mathbb C\), independent of \(E\),
for which
\[
 \chi(E,x)=\frac{y(E)}{P(E,x)}-\frac{iP_x(E,x)}{2P(E,x)}
 \quad\hbox{satisfies}\quad
 -i\chi_x+\chi^2+u-E=0
\]
for every \((E,x)\in U\times I\), if and only if \(Q(E,x)\) is independent
of \(E\). In that case \(u(x)=Q(E,x)\), equivalently
\begin{equation}
 R(E)=(E-u(x))P(E,x)^2+\tfrac12P(E,x)P_{xx}(E,x)
                    -\tfrac14P_x(E,x)^2.
 \label{eq:entire-identity}
\end{equation}
For \(x_0\in I\), the nonzero function
\[
 \psi(E,x)=\exp\!\left(\int_{x_0}^{x}
       \left(\frac{i y(E)}{P(E,s)}+\frac{P_x(E,s)}{2P(E,s)}\right)\,ds\right)
\]
satisfies \(-\psi_{xx}+u\psi=E\psi\) and \(\psi(E,x_0)=1\).
\end{proposition}
\begin{proof}
Put \(a=y/P\) and \(b=-P_x/(2P)\), so that \(\chi=a+ib\).
Because \(y\) is independent of \(x\), \(a_x=2ab\). Also
\(b_x=-P_{xx}/(2P)+P_x^2/(2P^2)\). Thus
\[
 -i\chi_x+\chi^2
   =a^2-b^2+b_x+i(2ab-a_x)
   =\frac{R}{P^2}-\frac{P_{xx}}{2P}+\frac{P_x^2}{4P^2}.
\]
The residual is exactly \(u-Q\). This proves both implications and
\eqref{eq:entire-identity}. The exponential is defined without selecting a
square root of \(P\). Its logarithmic derivative is \(i\chi\), so
\(\psi_{xx}/\psi=i\chi_x-\chi^2\); the differential equation follows.
\end{proof}

Equation \eqref{eq:entire-identity} is the classical squared-eigenfunction
identity; compare \citet[Eq.~(2.10)]{gesztesy1999infinite}. The proof above is
included to fix the signs in the notebook's convention, not to claim novelty.

\begin{proposition}[Invariance under a common factor \ours]
Suppose the compatibility criterion holds. For a holomorphic function
\(h(E)\ne0\) independent of \(x\), set
\(\widetilde P=hP\), \(\widetilde R=h^2R\), and \(\widetilde y=hy\).
Then \(\widetilde\chi=\chi\), \(\widetilde Q=Q\), and the potential is unchanged.
\end{proposition}
\begin{proof}
In the quotients \(\widetilde y/\widetilde P\),
\(\widetilde P_x/\widetilde P\), \(\widetilde P_{xx}/\widetilde P\), and
\(\widetilde R/\widetilde P^2\), all factors \(h\) cancel. This proves each
assertion. The nonvanishing hypothesis prevents spurious domain exclusions;
common zeros would require an additional removability discussion.
\end{proof}

\begin{example}[Two entire pairs of order one \ours]
The pair \(P(E,x)=R(E)=e^E\) consists of nonvanishing entire functions of
order exactly one, constant in \(x\). It gives
\(Q(E,x)=E-e^{-E}\), which is not independent of \(E\); therefore it defines
no potential through the notebook's displayed formula. In contrast,
\(P(E,x)=e^E\), \(R(E)=Ee^{2E}\) gives \(Q=0\), the free potential. Here
\(y=e^E\sqrt E\) can be chosen on any simply connected region avoiding zero,
and \(\psi=e^{i\sqrt E(x-x_0)}\). Both entire pairs meet the stated growth
condition. In the second example the exponential factors cancel: no infinite
spectral genus is created. More generally multiplying finite polynomial data
by \(h=e^{E^m}\), \(m\ge1\), changes growth without changing the potential.
Thus growth order alone neither guarantees compatibility nor specifies a
new class of potentials.
\end{example}

\paragraph{Independent second check.}
For any two solutions \(f,g\) of \(v''=(u-E)v\), put \(P=fg\) and
\(W=fg'-f'g\). Direct differentiation, independent of the Riccati substitution,
gives
\[
 (E-u)P^2+\tfrac12PP''-\tfrac14(P')^2=-\tfrac14W^2,
 \qquad W'=0.
\]
This rederives the same identity with \(R=-W^2/4\). The exact-polynomial
program \texttt{code/check\_spectral\_compatibility.py} verifies both derivations
and common-factor invariance over the rational numbers. A sign-sensitive
nonconstant check uses \(P=E+x^{-2}\), \(R=E^3\), \(u=2x^{-2}\) on \(x>0\)
away from zeros of \(P\); symbolic expansion and sixteen rational substitutions
give the same potential. \oursComputed{} The complete output is included. All
checks passed; the proof, rather than the finite substitutions, establishes the
general assertions.

\paragraph{Remaining scope.}
The local criterion does not characterize global real, nonsingular,
self-adjoint or almost-periodic realizations. A definite further problem must
choose the spectral variable (\(E\) or \(\sqrt E\)), remove common factors,
specify zero/divisor and reality conditions, and control high-energy behaviour
and continuation through zeros of \(P\). A further bounded research run could
test a chosen normalized family against those conditions. The present evidence
does not select such a family, so no new open conjecture is entered in the
problem index and no claim is made that the historical programme remains open
in every precise interpretation.
